% ID: FIS-FLU-PID-003
% Titolo: Ordinare pressioni a diverse profondita
% Disciplina: Fisica
% Area: Fluidostatica
% Argomento: Pressione idrostatica
% Sottoargomento: Confronto tra pressioni in un liquido
% Classe: Terza liceo artistico
% Difficolta: 2
% Tipo: quesito_teorico
% Risultato: p_A = p_B < p_C = p_D = p_E
% Tempo_stimato: 6 min
% Competenze: interpretazione fisica, confronto qualitativo, uso della legge di Stevino
% Prerequisiti: pressione idrostatica, profondita, densita
% Tag: fisica, fluidostatica, pressione_idrostatica, profondita, legge_stevino
% Fonte: verifica_fisica_fluidostatica_anonimizzata
% Autore: Riccardo Carli
% Licenza: CC-BY-SA-4.0

\begin{esercizio}
In un recipiente pieno d'acqua sono segnati cinque punti: $A$ e $B$ si trovano
alla stessa profondita, $C$ si trova piu in basso di $A$ e $B$, mentre $D$ ed
$E$ si trovano alla stessa profondita di $C$ ma in posizioni orizzontali
diverse.

Ordina le pressioni nei punti $A$, $B$, $C$, $D$, $E$ in ordine crescente e
motiva la risposta.
\end{esercizio}

\begin{soluzione}
In un liquido fermo la pressione idrostatica dipende solo dalla profondita:
\[
p=p_0+\rho gh.
\]
Non dipende dalla posizione orizzontale del punto.

I punti \(A\) e \(B\) sono alla stessa profondita, quindi hanno la stessa
pressione:
\[
p_A=p_B.
\]

I punti \(C\), \(D\) ed \(E\) sono piu in basso e alla stessa profondita tra
loro, quindi:
\[
p_C=p_D=p_E.
\]

Poiche \(C\), \(D\) ed \(E\) sono piu profondi di \(A\) e \(B\), la loro
pressione e maggiore. L'ordine crescente e:
\[
p_A=p_B < p_C=p_D=p_E.
\]
\end{soluzione}
